\documentclass[12pt]{article} % Clase de documento: artículo y tamaño de letra
\usepackage[utf8]{inputenc} % Escritura en castellano con acentos
\usepackage[T1]{fontenc} % Escritura en castellano con acentos
\usepackage{calligra} 
\usepackage{pslatex}  % Fuente de letras
\usepackage{graphicx}
\usepackage[margin=2.5cm]{geometry}
\usepackage{txfonts}
\usepackage{xcolor}
\usepackage{fancyhdr}
\usepackage{graphicx}
\usepackage{float}
\usepackage{listings}
\usepackage{multicol}
\usepackage{color}
\usepackage{times}
\usepackage{tikz}
\usepackage{verbatim}
\usepackage{mdwlist}
\usepackage{array}
\usepackage{caption}

\usetikzlibrary{chains,fit,shapes,arrows,calc,shapes,decorations.pathreplacing}


\pagestyle{fancy}
\headheight=50pt %para cambiar el tamaño del encabezado
\fancyhead[L]    %la "L" indica a la izquierda
{	
 \begin{minipage}{3.3cm}
  \includegraphics[width=1.0\textwidth]{Logo-UTN-BA.jpeg}
%  \includegraphics[width=0.6\textwidth]{UTNlogo.eps}
 \end{minipage}	
 \begin{minipage}{7.7cm}
   \normalsize
   {
     \textsf
     {
       \calligra{Universidad Tecnológica Nacional\\ Facultad Regional Buenos Aires \\ Departamento de {Ingeniería} {Electrónica}} 
     }
   }
 \end{minipage}
}

\fancyhead[R] %la "R" indica a la derecha
{
  \begin{minipage}{4.0cm}
   \small
   {
     \emph{\textbf{Informática I}} \\ \emph{\today} \\ \emph{Ejemplos de parcial}\\ \emph{Curso r1092}  
   }
  \end{minipage}}


\begin{document} % Inicio del documento
\newpage
%Cambia de página, el texto después de este comando aparecerá en la siguiente página en adelante.


\noindent
  \begin{center}
   \begin{tabular}{| c | c | c |}
    \hline
     Apellido y Nombres \hspace{8cm} &  Legajo & {Calificación} \\ \hline 
      &	& \\ \hline
   \end{tabular}	
  \end{center}

\noindent
Todas las funciones pedidas deben ser colocadas en el archivo {\bf primerParcialFunciones.c} y debe estar acompañado por su correspondiente {\bf punto h}.\\
Además se debe escribir un programa de testeo para cada función requerida. El nombre de esos archivos es {\bf primerParcialNumeroEjercicio.c}.\\
\\
%%%%%%%%%%%%%%%%%%%%%%%%%%%%%%%%%%%%%%%%%%%%%%%%%%%%%%%%%%%%%%%%%%%%%%%%%%%%%%%%%%%%%%%%%%%%%%%%%%%%%%%%%%%%% Comienzo del tema del examen
%%%%%%%%%%%%%%%%%%%%%%%%%%%%%%%%%%%%%%%%%%%%%%%%%%%%%%%%%%%%%%%%%%%%%%%%%%%%%%%%%%%%%%%%%%%%%%%%%%%%%%%%%%%%
\lstset{
	frame=Ltb,
	framerule=0pt,
	aboveskip=0.5cm,
	framextopmargin=3pt,
	framexbottommargin=3pt,
	framexleftmargin=0.4cm,
	framesep=0pt,
	rulesep=.4pt,
	backgroundcolor=\color{gray97},
	rulesepcolor=\color{black},
% 	language=C,
	captionpos=b,
	tabsize=3,
	frame=lines,
	keywordstyle=\color{blue},
	commentstyle=\color{Gray},
	stringstyle=\color{red},
	numbers=left,
	numberstyle=\tiny,
	numbersep=5pt,
	breaklines=true,
	showstringspaces=false,
	basicstyle=\ttfamily\scriptsize,
	emph={label},
	framerule=0pt,
}

Implemente una función que contenga el algoritmo de cifrado {\bf ROT13}. \\
El algoritmo consiste en sustituir cada letra por una que se encuentra trece posiciones por delante. Por ejemplo la A se reemplaza por la N, la B por la O y así sucesivamente. Para las últimas trece letras la secuencia se invierte. A continuación se muestra la tabla de las equivalencias entre las letras.
%\vspace{0.6cm}
\noindent
  \begin{center}
   \begin{tabular}{| c | c | c | c | c | c | c | c | c | c | c | c | c |}
    \hline
     A & B & C & D & E & F & G & H & I & J & K & L & M \\ \hline 
     N & O & P & Q & R & S & T & U & V & W & X & Y & Z \\ \hline 
   \end{tabular}	
  \end{center}
Ejemplo:
\noindent
\vspace{-0.6cm}
  \begin{center}
   \begin{tabular}{| c | c | c | c | c | c | c | c | c | c | c |}
    \hline
     Texto sin cifrar & H & O & L & A &   & M & U & N & D & O  \\ \hline 
     Texto cifrado    & U & B & Y & N &   & Z & H & A & Q & B  \\ \hline 
   \end{tabular}	
  \end{center}
Prototipo de la función:
\begin{center}
{\bf int rot\_13 (char *dataPtr)}
\end{center}
Donde {\bf dataPtr} es el puntero al mensaje a cifrar (string) y donde se colocará el mensaje cifrado y devuelve:
\begin{itemize}
\item {\bf -1} Si el mensaje contiene un caracter {\bf distinto} de una {\bf letra mayúscula} o un {\bf espacio}.
\item Un {\bf número positivo} indicando la cantidad de caracteres convertidos sin tener en 
cuenta los espacios ni el {\bf $\setminus$0}
\end{itemize}




 
\end{document} % Fin del documento.
 

  
